\section{Chapter plans}\label{sec:chapters}

Each chapter follows one template: problem, formulation, algorithm, implementation walkthrough with a tested listing of 15 to 40 lines, verification with measured numbers, failure modes stated as measured, and three to five exercises that point to a test or script. The plans below fix, for each chapter, the focus, the mathematics, the algorithm, the implementation mapping, the verification, the figures and tables, and the exclusions. ``Exists'' means the item is already in the repository; ``to produce'' means it must be created before the chapter is written.

\newcommand{\chapterplan}[4]{%
\subsection{Chapter #1: #2}\label{ch:#3}%
\textit{#4}\par\smallskip}

\subsection*{Part I: architectural foundations}

\chapterplan{1}{The design of a decoupled scientific engine}{arch}{Part I; 20 pages}
\begin{description}[style=nextline,leftmargin=1em]
\item[Focus] Separating materials, mesh data and element formulations from the solution drivers, so that a new element, material or solver is added by a registry entry or a subclass, not by editing existing code.
\item[Mathematics] Mapping every problem to the plain-array system $\bm{M}\ddot{\bm u}+\bm{C}\dot{\bm u}+\bm{K}\bm u=\bm F$; element interface as the pair stiffness and mass (and force and tangent for nonlinear elements).
\item[Algorithm] The flow material $\to$ mesh $\to$ element $\to$ system $\to$ solver $\to$ post-processing; how \code{solver.py} only calls \code{element.stiffness()} and reads \code{dofs_per_node}; backend dispatch between NumPy and PyTorch.
\item[Implementation] \anchor{fea_engine.solver.FESystem}, \anchor{fea_engine.elements.ELEMENT_REGISTRY}, \anchor{fea_engine.backend_dispatch}; the rule that \rom{} never imports \fea{} (only its tests do) and works on plain arrays.
\item[Verification] Registry conformance (every registered element passes the same interface test); overhead of dispatch against a direct call over increasing problem size. \todo{dispatch-overhead script and table.}
\item[Figures and tables] Architecture diagram in TikZ \todo{figure}; package map from the repository README \have{}; table of the six extension points (material, geometry, element, load, solution strategy, backend).
\item[Exclusions] Wrapper scripts, graphical interfaces, platform-specific build hooks.
\end{description}

\chapterplan{2}{Mesh topologies and structural domain partitioning}{mesh}{Part I; 22 pages}
\begin{description}[style=nextline,leftmargin=1em]
\item[Focus] Structured meshing for lines, rectangles and boxes; graded meshes; multi-block domains; holes and notches; and translation to a common exchange format.
\item[Mathematics] Node-connectivity arrays, coordinate transformations, grading functions $x_i = L\,(i/n)^p$, and block interfaces as shared node lists.
\item[Algorithm] Generators for 1-D, 2-D and 3-D grids; hole-in-rectangle and notch meshes; the single entry point \code{geometry.build_system}.
\item[Implementation] \anchor{fea_engine.mesh.rectangle_mesh}, \anchor{fea_engine.grading}, \anchor{fea_engine.geometry.build_system}, \anchor{fea_engine.mesh_io}; \code{MultiBlockMesh} for mixed element types.
\item[Verification] Index-mapping check of the exported connectivity against \code{meshio} (optional extra); element-quality and area/volume sum checks. \todo{meshio round-trip script.}
\item[Figures and tables] Hole and notch grading figures \todo{figure from the grading tests}; table of generators and the element each one targets.
\item[Design lesson] Why Gmsh support was removed, and what that did to the contract between mesh and element.
\item[Exclusions] Unstructured generators, Delaunay triangulation, manual input-file editing.
\end{description}

\subsection*{Part II: full-order physics (\fea{})}

\chapterplan{3}{Linear continua and structural elements}{elem}{Part II; 34 pages}
\begin{description}[style=nextline,leftmargin=1em]
\item[Focus] From weak forms to sparse element matrices, for scalar-field, plane, solid, beam, plate and shell elements.
\item[Mathematics] Variational equation for the 4-node Laplace element; plane stress and strain; mean-dilatation B-bar for the 8-node hexahedron~\cite{hughes1980bbar}; Mindlin plate; MITC shell~\cite{dvorkin1984}; Euler--Bernoulli, corotational~\cite{crisfield1990corot} and Reissner~\cite{simo1985beam} beams.
\item[Algorithm] Gauss integration loops; consistent against lumped mass; reduced integration and hourglass stabilisation; source vectors.
\item[Implementation] \anchor{fea_engine.elements.scalar.Quad4Laplace}, \anchor{fea_engine.elements.solids}: \code{Quad4PlaneStress}, \code{Quad8PlaneStress}, \code{Tri3PlaneStress}, \code{Tri6PlaneStress}, \code{Hex8Solid3D}, \code{Hex8SolidBbar}, \code{Hex20Solid3D}, \code{Tet4Solid3D}, \code{Tet10Solid3D}; \anchor{fea_engine.elements.beams}, \code{beams3d}, \code{plates}, \code{shells}, \code{shells_director}.
\item[Verification] Patch tests (paper table \code{tab_patch.tex}) \have{}; rigid-body and null-space tests; cantilever against Euler--Bernoulli (0.29\% at 128 elements in \code{fea_01}) \have{}; shell verification figure \have{}; membrane eigenvalues for \code{Quad4Laplace} \have{}.
\item[Figures and tables] \code{fea_01}, \code{fea_05} and the shell verification figure \have{}; table of elements with DOFs per node and integration rule \have{} (README module reference); locking comparison of \code{Hex8Solid3D}, \code{Hex8SolidBbar} and \code{Hex20Solid3D} \todo{figure from the locking test}.
\item[Exclusions] Structural-engineering case studies, fatigue, failure theory, basic shape-function derivations.
\end{description}

\chapterplan{4}{Dynamics: modal, transient, harmonic and random response}{dyn}{Part II; 28 pages; new}
\begin{description}[style=nextline,leftmargin=1em]
\item[Focus] The dynamic solution strategies that the reduced-order chapters reduce and that the contact chapter extends.
\item[Mathematics] Generalised eigenproblem; Newmark family~\cite{newmark1959} and its stability; critical time step $\Delta t_c \le 2/\omega_{\max}$; modal superposition with damping ratios; harmonic response; random vibration with power spectral densities.
\item[Algorithm] Implicit Newmark--$\beta$ with one factorisation; explicit central differences with a lumped mass matrix; mode truncation; frequency sweeps.
\item[Implementation] \anchor{fea_engine.solver.FESystem.solve_modal}, \code{solve_transient_implicit}, \code{solve_transient_explicit}, \code{critical_timestep}, \code{solve_modal_superposition}, \code{solve_harmonic}, \code{solve_random_vibration}; Rayleigh damping in \code{damping.py}.
\item[Verification] Mode 1 of the 3-D cantilever 21.107~Hz against Euler--Bernoulli 20.955~Hz \have{}; step-response peak over static value 1.927 with 1\% damping \have{}; Newmark order-of-accuracy table \todo{script}; explicit stability-limit test \todo{script}.
\item[Figures and tables] \code{fea_05} mode shapes \have{}; \code{fea_08} transient response (animated GIF) \have{}; accuracy and stability table \todo{table}.
\item[Exclusions] Wave-propagation case studies, chaotic response, structural-health monitoring.
\end{description}

\chapterplan{5}{Material and geometric nonlinearity}{nonlin}{Part II; 36 pages}
\begin{description}[style=nextline,leftmargin=1em]
\item[Focus] Total-Lagrangian kinematics, hyperelasticity and J2 plasticity, and the solution paths that follow equilibrium through limit points.
\item[Mathematics] Green--Lagrange strain $\bm E=\tfrac12(\bm F^T\bm F-\bm I)$, second Piola--Kirchhoff stress, Neo-Hookean energy, return mapping with isotropic and kinematic hardening~\cite{simo1998inelasticity}, consistent tangent~\cite{simo1985tangent}; arc-length~\cite{riks1979,crisfield1981} and Koiter--Newton~\cite{liang2014koiternewton} continuation.
\item[Algorithm] Newton--Raphson with load stepping, displacement control, arc-length, Koiter--Newton, line search and the shared \code{NewtonOptions}.
\item[Implementation] \anchor{fea_engine.elements.nonlinear_solids}: \code{Tet4NeoHookean}, \code{Quad4PlaneStressTL}, \code{Tet10SolidTL}, \code{Hex8PlasticJ2}, \code{Hex8PlasticJ2Kinematic}, \code{Quad4PlasticJ2PlaneStress}; \anchor{fea_engine.elements.hex_tl}; \code{TrussTL2D}; the separate \code{state} and \code{iter_state} that keep path dependence; \anchor{fea_engine.nonlinear_solver}.
\item[Verification] Elastica tip displacement (paper table \code{tab_elastica.tex}) \have{}; snap-through with arc-length \have{}; plasticity load--unload loop \have{}; finite-difference tangent checks for the total-Lagrangian elements \have{}.
\item[Figures and tables] \code{fea_02}, \code{fea_04}; snap-through and uniaxial J2 loop figures in \code{fea_engine/examples} \have{}; driver-choice table \have{} (README).
\item[Exclusions] Viscoelasticity, crack propagation, damage.
\end{description}

\chapterplan{6}{Computational acceleration: vectorised assembly and batched kernels}{batch}{Part II; 30 pages}
\begin{description}[style=nextline,leftmargin=1em]
\item[Focus] Removing the Python loop over elements, keeping the tangent sparse, and knowing where this helps and where it does not.
\item[Mathematics] Element kernels as contractions over an element axis and a Gauss-point axis; the element tangent $\bm K_e=\sum_g w_g J_g\,\bm B_g^T\bm D\bm B_g$ folded\linebreak into one product of contraction length $6G$; the CSR position map and \code{bincount} assembly.
\item[Algorithm] Stacked coordinates $(E,n,d)$ and displacements $(E,n_{dof})$; memory-bounded chunking; per-block fallback to the loop; sparse LU in the Newton drivers with optional factorisation reuse.
\item[Implementation] \anchor{fea_engine.batched_kernels}, \anchor{fea_engine.vectorized_assembly}, \anchor{fea_engine.batched_kernels_torch}, \anchor{fea_engine.contact_batch}; keywords \code{vectorized}, \code{chunk_size}, \code{sparse}, \code{sparse_tangent}.
\item[Verification] Agreement with the loop to $10^{-12}$ in tests and $8.4\times10^{-15}$ at most in the measured cases; speedups 102x, 47x, 23x and 41x for force on \code{Quad4PlaneStressTL}, \code{Hex8SolidTL}, \code{Hex20SolidTL} and \code{Tet10SolidTL} (\code{paper/results/batched_kernels.json}) \have{}.
\item[Figures and tables] Speedup and agreement table \have{} (data), figure \todo{figure from the JSON}; chunking memory table \todo{from the foundation tests}; torch CPU/CUDA comparison \have{} (data).
\item[Honest limits] The 3-D gain is mainly the folded kernel, the sparse tangent and the cheaper solve; the GPU was not faster than torch on the CPU in the one case timed; $10^5$-element meshes were not tested.
\item[Exclusions] Multi-threading design, C++ extension bindings, CPU-cache tuning.
\end{description}

\chapterplan{7}{Unilateral contact mechanics and sparse active sets}{contact}{Part II; 36 pages}
\begin{description}[style=nextline,leftmargin=1em]
\item[Focus] Exact enforcement of frictionless contact with Lagrange multipliers instead of a penalty parameter.
\item[Mathematics] The discrete conditions
\[
\bm K\bm u=\bm f+\bm B^T\bm\lambda,\qquad \bm B\bm u-\bm c\ge 0,\qquad \bm\lambda\ge 0,\qquad \bm\lambda^T(\bm B\bm u-\bm c)=0,
\]
their equivalence to a linear complementarity problem and to a semismooth Newton method~\cite{hintermuller2002pdas}; the contact formulation in the literature~\cite{wriggers2006contact,wohlmuth2011contact}; BDF1 and BDF2 time integration with contact.
\item[Algorithm] Primal-dual active set with the routes \code{reduced} ($\bm B$ selects nodes), \code{saddle} (general sparse $\bm B$) and \code{condensed} (dense reference); warm start; failure reports with \code{reason}, \code{message} and \code{suggestion}; BDF2 contact dynamics.
\item[Implementation] \anchor{fea_engine.contact_solvers.solve_contact_static_sparse}, \anchor{fea_engine.contact_solvers.solve_contact_dynamics}, \code{ContactResult}, \code{ContactDynamicsInfo}; the scalar element \anchor{fea_engine.elements.scalar.Quad4Laplace}; penalty elements \code{GapContactPenalty} and \code{NodeToSegmentContact2D} as the contrast.
\item[Verification] Routes against dense active-set and Lemke solutions to $1.3\times10^{-15}$ (361 to 3,481 unknowns) and $7.4\times10^{-15}$ between sparse routes at 39,601 unknowns \have{}; 39,601 unknowns with 9,425 active constraints in 4.5~s (reduced) and 7.6~s (saddle) \have{}; BDF2 observed order 1.96 and 2.03 against 0.92 for the default scheme \have{}; energy history of the bar impact \have{}.
\item[Figures and tables] \code{fea_03}, \code{fea_09}, \code{fea_10} \have{}; route table; failure-mode table built from the failure-report tests \have{}.
\item[Honest limits] BDF2 is not contact-stabilised (energy rise 0.6\% to 0.72\% on impact, wall-force spikes); no convergence guarantee when $\bm K$ is not an M-matrix; friction, Newmark contact and the dense route at 39,601 unknowns were not run.
\item[Exclusions] Multi-body friction laws, adhesion, smoothing of flexible-to-flexible impact.
\end{description}

\subsection*{Part III: reduced-order modelling (\rom{})}

\chapterplan{8}{Projection-based reduction: POD, Galerkin and affine sweeps}{pod}{Part III; 28 pages}
\begin{description}[style=nextline,leftmargin=1em]
\item[Focus] Compressing a large linear model onto a few modes and sweeping parameters at the cost of a small solve.
\item[Mathematics] Proper orthogonal decomposition via the singular value decomposition~\cite{berkooz1993pod}, the energy (mass) inner product $\langle x,y\rangle_M=x^T\bm M y$, Galerkin projection $\bm K_r=\bm V^T\bm K\bm V$, affine parameter dependence $\bm K(\mu)=\sum_q\theta_q(\mu)\bm K_q$~\cite{hesthaven2016rb}.
\item[Algorithm] Snapshot sampling (including optimal Latin hypercube designs), mass-weighted POD, offline--online split, a posteriori reduced-versus-full comparison on unseen parameters.
\item[Implementation] \anchor{rom_engine.pod}, \anchor{rom_engine.galerkin}, \anchor{rom_engine.affine}, \anchor{rom_engine.parametric_modes}, \anchor{rom_engine.sampling}, \anchor{rom_engine.metrics}.
\item[Verification] Singular-value decay and projection error on unseen parameters; 500-point affine sweep, full order 1653.23~ms against reduced 7.45~ms (221.8x), maximum relative error $1.51\times10^{-3}$ \have{}; 8 snapshots, 6-mode basis, error about $7\times10^{-6}$ \have{}.
\item[Figures and tables] \code{rom_01}, \code{rom_02}, paper figure e1 \have{}; table of the offline and online cost split \todo{table from e1 data}.
\item[Exclusions] Non-projection reduction, regression models, purely statistical surrogates.
\end{description}

\chapterplan{9}{Systems-and-control reduction}{bt}{Part III; 26 pages}
\begin{description}[style=nextline,leftmargin=1em]
\item[Focus] Compact input--output models of structural systems, with error bounds.
\item[Mathematics] First- and second-order state space; Lyapunov equations $\bm A\bm P+\bm P\bm A^T+\bm B\bm B^T=0$, controllability and observability Gramians, Hankel singular values and the bound $\|G-G_r\|_\infty\le 2\sum_{k>r}\sigma_k$~\cite{moore1981,glover1984}; singular perturbation approximation~\cite{liu1989spa}; frequency-weighted Gramians~\cite{enns1984}; Krylov moment matching~\cite{antoulas2005} and second-order Arnoldi~\cite{bai2005soar}.
\item[Algorithm] Square-root balanced truncation, SPA, frequency-weighted truncation, one- and two-sided Krylov projection, SOAR.
\item[Implementation] \anchor{rom_engine.state_space}, \anchor{rom_engine.balanced_truncation}, \anchor{rom_engine.hankel_norm}, \anchor{rom_engine.krylov}, \anchor{rom_engine.soar}, \anchor{rom_engine.passivity}, \anchor{rom_engine.slow_fast}.
\item[Verification] Actual error below the certified bound over a frequency sweep \have{} (\code{rom_04}); Krylov one-sided against two-sided mean error 0.153 against 0.0115 \have{}; passivity check \have{}.
\item[Figures and tables] \code{rom_04}; comparison table of the reduction methods \have{} (user guide); comparison with the method comparison of Besselink et al.~\cite{besselink2013comparison}.
\item[Exclusions] Controller design, PID tuning, feedback software.
\end{description}

\chapterplan{10}{Frequency-domain ROMs and certified error bounds}{freq}{Part III; 26 pages}
\begin{description}[style=nextline,leftmargin=1em]
\item[Focus] Fast steady-state frequency responses with a rigorous error estimate and an adaptive training loop.
\item[Mathematics] Successive constraint method for the coercivity lower bound~\cite{huynh2007scm}, residual-based a posteriori error estimates, random-vibration power spectral density.
\item[Algorithm] Greedy frequency training driven by the residual indicator; SCM bound by linear programming; PSD propagation through the reduced transfer function.
\item[Implementation] \anchor{rom_engine.frequency}, \anchor{rom_engine.greedy}, \anchor{rom_engine.scm}, \anchor{rom_engine.scm_lp}, \anchor{rom_engine.random_vibration}.
\item[Verification] Greedy against uniform training at equal budget \have{} (\code{rom_03}); residual curve against the certified bound over a dense sweep \have{}.
\item[Figures and tables] \code{rom_03}; certified-bound plot \todo{figure from the certified-bound example}.
\item[Exclusions] Transient wave modulation, chaotic response.
\end{description}

\chapterplan{11}{Substructuring and component reduction}{cms}{Part III; 22 pages; new}
\begin{description}[style=nextline,leftmargin=1em]
\item[Focus] Reducing each component of an assembly and coupling the components through reduced interfaces.
\item[Mathematics] Static condensation~\cite{guyan1965}, Craig--Bampton fixed-interface modes~\cite{craig1968}, interface reduction.
\item[Algorithm] Component mode synthesis, component POD, interface mode construction (joint, Gram--Schmidt).
\item[Implementation] \anchor{rom_engine.cms}, \anchor{rom_engine.component_pod}, \anchor{rom_engine.interface_reduction}.
\item[Verification] Reduced assembly against the full assembly (paper experiment e3, table \code{tab_cms.tex}) \have{}.
\item[Figures and tables] Paper figure e3 \have{}; interface-size against accuracy table \todo{from the e3 data}.
\item[Exclusions] Multi-scale domain decomposition algorithms.
\end{description}

\chapterplan{12}{Non-intrusive identification and nonlinear ROMs}{nonint}{Part III; 32 pages}
\begin{description}[style=nextline,leftmargin=1em]
\item[Focus] Reduced models when the full-order matrices are not available or the nonlinearity prevents direct projection.
\item[Mathematics] Loewner interpolation from frequency data~\cite{mayo2007loewner}; implicit condensation and expansion with polynomial stiffness terms~\cite{hollkamp2008ice}; multi-fidelity surrogate force models with radial basis functions~\cite{he2023mfsnlrom}; nonlinear normal modes and backbone curves~\cite{kerschen2009nnm}; harmonic balance and continuation~\cite{keller1977}; quadratic manifolds.
\item[Algorithm] Loewner pencil and physical-mode screening; operator inference; force-coefficient identification; polynomial modal ROM; backbone tracing.
\item[Implementation] \anchor{rom_engine.loewner}, \anchor{rom_engine.screening}, \anchor{rom_engine.operator_inference}, \anchor{rom_engine.nonlinear_rom}, \anchor{rom_engine.nnm}, \anchor{rom_engine.quadratic_manifold}, \anchor{rom_engine.forced_response_hbm}.
\item[Verification] Held-out load histories not used in sampling; 2-mode nonlinear modal ROM against a 51-DOF beam, maximum error 3.3\% \have{} (\code{rom_05}); the MFS-NLROM, ICE and NonLin-HyROM reproductions \have{} (\code{paper/reproductions}).
\item[Figures and tables] \code{rom_05}; \code{repro_mfs_backbone}, \code{repro_mfs_sweep}, \code{repro_ice_1mode}, \code{repro_hyrom_case1_backbone}, \code{repro_hyrom_case2_backbone} \have{}; surrogate comparison table (paper \code{tab_surrogates.tex}) \have{}.
\item[Note on neural models] One paragraph on \code{neural_operator.py} and \code{parameterized_latent_ode.py} as optional tools, (decided: one paragraph, and no neural-network chapter).
\item[Exclusions] Deep-network architectures, general machine-learning tutorials, reinforcement learning.
\end{description}

\chapterplan{13}{Hyper-reduction and contact reduced-order models}{hyper}{Part III; 34 pages}
\begin{description}[style=nextline,leftmargin=1em]
\item[Focus] Removing the last full-order cost, the internal and contact force evaluation, and keeping contact conditions valid in the reduced space.
\item[Mathematics] Energy-conserving sampling and weighting~\cite{farhat2015ecsw}; a non-negative basis for contact forces and the reduced contact problem~\cite{balajewicz2016contact}; the hybrid reduced-integration-domain method~\cite{fauque2018hybrid}; static condensation of the linear zone.
\item[Algorithm] ECSW training matrix and non-negative least squares; contact-aware weights with the contact zone held at weight one; reduced LCP; reduced BDF2 contact dynamics with warm-started active sets and a KKT guard; greedy snapshot selection by the violation of the contact conditions with a cached indicator.
\item[Implementation] \anchor{rom_engine.hyper_reduction}, \anchor{rom_engine.contact_dual_basis.ContactROM}, \anchor{rom_engine.contact_greedy}, \anchor{rom_engine.contact_hyper_reduction}, \anchor{rom_engine.contact_rid}, \anchor{rom_engine.domain_partition}, \anchor{rom_engine.complementarity}.
\item[Verification] Complete bases reproduce the full-order trajectory to $3.3\times10^{-11}$; error 0.81\%, 0.20\%, 0.023\% for $p=p_\lambda=10,30,40$ with reduced KKT residual at most $2.7\times10^{-18}$ \have{}; greedy against random selection, 0.82\% against 1.63\% (median) at 8 snapshots \have{}; ECSW training matrix 92x faster and the same 13 elements selected \have{}.
\item[Figures and tables] \code{rom_07}, paper figure e4 \have{}; indicator timing table (0.04 to 0.45~ms per call up to 40,000 constraints) \have{}.
\item[Honest limits] Parametric prediction in the dynamic case and friction were not tested; the \code{rel\_tol} stopping rule is an author's reading of line 15 of Algorithm~2 of the cited method.
\item[Exclusions] Global contact re-meshing, multi-scale domain decomposition.
\end{description}

\subsection*{Part IV: verification}

\chapterplan{14}{Verification, reproduction studies and a parametric pipeline}{verif}{Part IV; 30 pages}
\begin{description}[style=nextline,leftmargin=1em]
\item[Focus] Treating verification as part of the architecture, and showing the suite on published problems and on a complete parametric workflow.
\item[Mathematics] Order-of-accuracy estimation, energy balances, relative-error formulas normalised by the reference maximum.
\item[Algorithm] Cross-validation sweeps over element families; cross-configuration validation of reduced models; offline construction followed by online queries on unseen inputs.
\item[Implementation] \anchor{rom_engine.validation}, \anchor{rom_engine.metrics}, \anchor{rom_engine.cross_configuration_validation}, \anchor{rom_engine.dataset_diagnostics}, \anchor{rom_engine.ensemble_uq}; the scripts of \code{paper/experiments} (e0 to e10).
\item[Verification] Verification table against closed-form targets (paper experiment e5, \code{tab_patch.tex}, \code{tab_orders.tex}) \have{}; five reproduction reports \have{} (ICE, MFS-NLROM, NonLin-HyROM, fluid--structure modal identification, FD-PINN); end-to-end workflow with unseen loads (experiment e8) \have{}.
\item[Figures and tables] Paper figures e5, \code{repro_fsi_mac}, \code{repro_fsi_noise}, \code{repro_fdpinn_*} \have{}; summary table of the reproductions \todo{table assembled from the five reports}.
\item[Exclusions] Web dashboards, cloud deployment, sensor integration, a ``digital twin'' claim; only measured speedups are stated.
\end{description}

\subsection*{Appendix}

\chapterplan{A}{Differentiable assembly and autograd tangents (optional)}{diff}{Appendix; 12 pages}
\begin{description}[style=nextline,leftmargin=1em]
\item[Focus] Exact tangents and gradients through the solver when PyTorch is present.
\item[Mathematics] Implicit differentiation of a Newton solve, adjoint and forward modes, trainable additive corrections.
\item[Implementation] \anchor{fea_engine.differentiable}, \anchor{fea_engine.autograd_tangent}, \anchor{rom_engine.differentiable_correction}.
\item[Verification] Torch-free checks regenerated by the book script (FD floor of the tetrahedron tangent against a complex-step reference, Newton with a correction, adjoint gradient against central differences, metric identity) \have{}; autograd tangents against the other methods at the tolerances asserted in the PyTorch-gated tests, quoted and not run \have{}. The figure $1.25\times10^{-16}$ that the plan once listed here is the agreement of a NumPy mirror of the shell force with the production force, not an autograd result.
\item[Exclusions] Training of neural models and optimisation frameworks.
\end{description}
